\paragraph{About this part.} Part~II asked what the book predicts. Part~III is about the infrastructure needed to
turn a prediction into a tested trading decision. It sets out what a simulator must capture for execution
questions---queue position, adverse selection and latency (Section~\ref{sec:sim-needs}); surveys the order-book
simulators available and compares them with the Kaspar system, including how an order-book model is added to Kaspar
as an actor that publishes its predictions to a strategy (Section~\ref{sec:simulators}); explains why order-book
models are slow at inference and what hardware can accelerate (Section~\ref{sec:hardware}); and closes with open
problems and the survey's conclusion (Sections~\ref{sec:agenda} and~\ref{sec:conclusion}).

% ================================================================
\section{Mechanics a simulator must capture}
\label{sec:sim-needs}

A simulator for execution research has to reproduce the mechanics that decide whether an order trades and what
the trade is worth. This section sets out those mechanics: the difference between crossing the spread and
waiting in the queue, how queue position determines whether a passive order is filled, why the fills it does get
are biased towards bad ones, and how latency reduces the value of a signal. Section~\ref{sec:simulators} then
compares simulators on how well they model each of them.

\subsection{Passive and aggressive execution}
\label{sec:passive}

An order either crosses the spread and trades at once (aggressive) or rests in the book and waits to be filled
(passive). This section writes the profit of each and splits the passive profit into a fill probability and a value
per fill.

Consider one unit traded on side $d$ at time $t$. For an aggressive order, which executes
immediately, a simplified realised profit is
\begin{equation}
  \Pi^{\text{agg}} = d\,\Delta m_{t,H} - C_{\text{spread}} - C_{\text{fees}} - C_{\text{impact}},
\end{equation}
where $C_{\text{spread}} = s_t/2$ is the half spread paid to cross, $C_{\text{fees}}$ the exchange
and clearing fees, and $C_{\text{impact}}$ the additional price concession from walking the book.
\inwords{profit from trading immediately is the price move in our favour---multiplying by $d$
flips the sign for a sell, so a falling price is good for a seller---minus what it cost to get in:
half the spread, the fees, and any extra paid for consuming more than the best level.}
A passive order placed at $t$ trades only if it is filled, at time $t_F$ and price
$p^{\text{fill}}$. Define its post-fill mark-out over horizon $h$ as
\begin{equation}
  R^{\text{fill}}_{h} = d\left(m_{t_F+h} - p^{\text{fill}}\right).
  \label{eq:markout}
\end{equation}
\inwords{once our order has traded, look at the mid $h$ seconds later and ask how far it is from
the price we got, counted positive when it is in our favour. A large positive mark-out means we
bought cheaply; a negative one means the market moved against us right after we traded.}
For a passive buy filled at the bid, the mark-out is the half spread at the fill time,
$s_{t_F}/2$, plus the mid move after the fill. The passive profit is
\begin{equation}
  \Pi^{\text{pass}} = \mathbb{1}_{\Fill}\left(R^{\text{fill}}_{h} + R_{\text{rebate}}\right),
\end{equation}
where $R_{\text{rebate}}$ is any liquidity-provision rebate net of fees.
\inwords{if the order is filled, the profit is its mark-out plus any rebate the exchange pays for
providing liquidity; if it is never filled, the indicator is zero and so is the profit.} The passive profit exists
only conditional on execution, which gives the decomposition
\begin{equation}
  \E[\Pi^{\text{pass}}] = \P(\Fill)\; \E[\Pi^{\text{pass}} \mid \Fill].
  \label{eq:decomp}
\end{equation}
\inwords{the average profit of a passive order equals the chance it is filled times the average
profit when it is. This is the identity $\E[Z\,\mathbb{1}_A] = \P(A)\,\E[Z \mid A]$ of
Section~\ref{sec:primer}, with $A$ the fill.}
Equation~\eqref{eq:decomp} separates two quantities that a directional forecast conflates: how
likely the order is to trade, and how good the trade is when it does. The rest of this section is
about the two factors. Market-making models~\citep{avellaneda2008,cartea2015} and optimal-execution
models~\citep{almgren2001} give the surrounding control problems; here we focus on the inputs
those problems need from the book.

\subsubsection{Signal value and execution value}
\label{sec:tradeability-framework}

The decomposition gives tradeability a precise meaning. A directional model estimates the
\emph{signal value}
\begin{equation}
  V^{\text{sig}}(X_t) = \E[\Delta m_{t,H} \mid X_t],
\end{equation}
the expected price move given what is known. A passive execution decision is worth its
\emph{execution value}
\begin{equation}
  V^{\text{exec}}(X_t, a_t) = \P(\Fill \mid X_t, a_t)\; \E\big[\Pi^{\text{pass}} \mid \Fill, X_t, a_t\big],
\end{equation}
which depends on the action through where the order is placed and therefore on its queue position.
\inwords{signal value asks ``which way will the price go?''; execution value asks ``if I act on that
view with this order, how likely am I to trade, and how good will the trade be when I do?'' A model
can be excellent at the first and useless at the second.}

The two can even point in opposite directions. A forecast of a falling price is most confident when
the book is already one-sided---a thin bid, heavy selling. That is exactly when a resting buy order is
most likely to be filled and least likely to be worth having: the fill and the adverse move arrive
together (Section~\ref{sec:queue-models}). Conditioning on the fill reverses the apparent value of the
signal. This is why the survey treats $\P(\Fill \mid \cdot)$ and $\E[\,\cdot \mid \Fill]$ as separate
modelling targets, and why evaluation must use the conditional quantities rather than unconditional
accuracy.

Who earns the execution value in practice is documented. A large high-frequency market maker studied
by \citet{menkveld2013} traded passively in about four of every five trades, earning the spread while
losing on its inventory; across high-frequency firms, differences in relative speed explain large
differences in trading performance~\citep{baron2019}.

\subsection{Queue position and fill probability}
\label{sec:queue-position}

A passive order is filled only when everything ahead of it in its queue has been traded or
cancelled. This section shows why that makes queue position a variable in its own right: two orders
at the same price, with the same forecast behind them, can have very different chances of being
filled.

\subsubsection{Residual queue position}

For a resting order, the displayed size at its price does not determine its fill probability. Let
$Q^{\text{ahead}}_t$ be the volume ahead of the order at the same price. Under price--time
priority the order cannot execute until that volume has been removed by executions or
cancellations. Two orders at the same price can therefore have very different values, and queue
position has been modelled explicitly as a source of value~\citep{moallemi2017}.

For a passive order with $q$ ahead of it, the probability of a fill within horizon $H$ is
$\P(\Fill \mid Q^{\text{ahead}}_t = q,\, X_t)$. It must be distinguished from
$\P(\Delta m_{t,H} > 0 \mid X_t)$. A model can predict direction excellently and still have poor
trading value if the predicted moves happen mainly when the passive order is unlikely to execute.

\begin{example}[Same forecast, different orders]
\label{ex:queue}
Two resting bids at the same price have 20 and 480 lots ahead of them. Net depletion runs at 40
lots per second. The naive time to reach the front is $0.5$ s for the first and $12$ s for the
second. Both carry the same $\P(\Delta m_{t,H} > 0 \mid X_t)$. Over a five-second horizon the first
almost surely fills and the second almost surely does not.
\end{example}

\subsection{Adverse selection}
\label{sec:adverse}

Being filled is not always good news. A passive order tends to be filled exactly when someone better
informed wants to trade against it, so the fills it gets are, on average, followed by price moves
against it. This section defines that cost precisely and shows how it depends on queue position.

\subsubsection{Conditioning on the fill}

Using the mark-out $R^{\text{fill}}_{h}$ of~\eqref{eq:markout}, define the conditional
adverse-selection component for an order that had $q$ ahead of it when placed as
\begin{equation}
  \mathrm{AS}(q, h) = -\E\left[R^{\text{fill}}_{h} - \tfrac{s_{t_F}}{2} \,\middle|\, \Fill,\, Q^{\text{ahead}}_t = q\right],
\end{equation}
the expected mid move against the order after it is filled, excluding the half spread it earns.
\inwords{among orders that started with $q$ lots ahead of them \emph{and were filled}, how far,
on average, did the price then move against us? It is a form of the winner's curse: in an auction,
the times your low bid wins are disproportionately the times when others knew something you did
not. A passive order is most likely to be filled just when an informed trader wants to trade the
other way.}
The conditioning on the fill is essential. Without it, the quantity measures the unconditional
drift of the market rather than the information in the event that the order actually traded.
Classical models make the mechanism explicit: a liquidity provider who trades with a better-informed
counterparty loses on average~\citep{glosten1985,kyle1985}.

\subsubsection{Queue position as an information variable}

Queue position affects both the probability of execution and the information content of
execution. \citet{moallemi2017} model both effects and find that a position at the front of the
queue is always worth more than one at the back: in their model adverse selection increases with queue
position, because an order at the end of a long queue tends to be filled only by a large aggressive
order, which is more likely to be informed. For some large-tick US stocks they estimate the front--back
value difference to be comparable to the bid--ask spread. The direction of the effect is therefore
supported by a model; its size is an empirical object that can depend on asset, tick regime,
spread, order-flow state, horizon and cancellation behaviour, and execution-aware benchmarks should
report it directly (Section~\ref{sec:agenda}).

\begin{example}[The decomposition at two positions]
\label{ex:as}
All values in ticks; fill window $H = 10$ s, mark-out horizon $h = 10$ s, one-tick spread, no
rebate. Here ``$\mid q$'' abbreviates $\mid Q^{\text{ahead}}_t = q$, and

\[
  \E[\Pi^{\text{pass}} \mid \Fill, q] = s_{t_F}/2 - \mathrm{AS}(q,h):
\]

\begin{center}\small
\begin{tabular}{lccccc}
\toprule
Position & $\P(\Fill \mid q)$ & $s_{t_F}/2$ & $-\mathrm{AS}(q,h)$ & $\E[\Pi^{\text{pass}} \mid \Fill, q]$ & $\E[\Pi^{\text{pass}} \mid q]$ \\
\midrule
Front & 0.80 & $+0.5$ & $-0.20$ & $+0.30$ & $+0.24$ \\
Back  & 0.25 & $+0.5$ & $-0.45$ & $+0.05$ & $+0.0125$ \\
\bottomrule
\end{tabular}
\end{center}
The table shows the arithmetic of~\eqref{eq:decomp}; the ordering of the two rows matches the
direction found by~\citet{moallemi2017}, with stylised magnitudes. A metric that looked
only at direction or only at fill rate would rank the two positions without seeing the product.
\end{example}

\subsection{Latency and signal decay}
\label{sec:latency}

A prediction at horizon $H$ is useful only if its information survives until the decision is
executed. This section splits the delay between seeing the data and acting on it into stages and shows how it
reduces the value of a signal. Write the total latency as
\begin{equation}
  T_{\text{lat}} = T_{\text{signal}} + T_{\text{decision}} + T_{\text{network}} + T_{\text{exchange}},
\end{equation}
the times taken to compute the signal, to decide on an action, to transmit the order, and for the
exchange to process it.
\inwords{total delay from seeing the data to the order being live at the exchange is the sum of
the stage delays, just as the latency of a processor pipeline is the sum of its stages.}
A signal whose value decays materially before $T_{\text{lat}}$ is not economically equivalent to
one of the same accuracy with longer persistence. The relevant object is not
$\E[\Delta m_{t,H} \mid X_t]$ but
\begin{equation}
  \E\left[\Pi \mid X_t,\, T_{\text{lat}},\, \text{execution policy}\right],
\end{equation}
the expected profit given the information, the latency, and the rule that maps forecasts to orders.
\inwords{the quantity that matters is not ``will the price go up?'' but ``how much do we expect to
make, given what we knew, how slow we are, and what we do with the forecast?''}
Trading firms compete on latency, and that competition has itself been studied as an economic
phenomenon~\citep{budish2015,hasbrouck2013}.

\begin{example}[Half-life against latency]
Let a signal's value decay exponentially with a half-life of 2 ms, so a fraction
$2^{-T_{\text{lat}}/2\,\mathrm{ms}}$ of it remains after a delay $T_{\text{lat}}$. At $100\,\mu$s total latency
$2^{-0.05} \approx 97\%$ of its value remains; at 5 ms only $2^{-2.5} \approx 18\%$ does. The same
forecast accuracy yields very different economics.
\end{example}

% ================================================================
\section{Order-book simulators}
\label{sec:simulators}
% ================================================================

An order-book simulator estimates what would have happened to a researcher's orders. Simulators differ in where the
market comes from and in how faithfully they model the mechanics of Section~\ref{sec:sim-needs}. This section first
shows why a naive backtest is not a simulator for passive orders and what a queue-exact replay requires; it then
surveys the simulators in use for research and compares them on those requirements, describes the Kaspar system in
detail, and shows how an order-book model is added to it.

\subsection{Naive backtests and queue-exact replay}
\label{sec:replay}

The questions of Section~\ref{sec:sim-needs}---will the order fill, and is the fill a good one---cannot be answered
by a backtest that assumes orders fill when the price touches them. This section shows how such a backtest misleads,
lists what a simulation must model to give trustworthy answers, and states what even an exact replay cannot see.

\subsubsection{Assumptions of a naive backtest}

A common backtest of a passive signal assumes that a limit order placed at the bid is filled when
the signal fires, or when the price touches the order's level:
\begin{lstlisting}
naive:   signal at t -> assume filled at bid -> P&L = m(t+H) - bid
\end{lstlisting}
This fills every order, including those the market would never have reached. The orders it wrongly
fills are exactly the ones where the price moved away before the queue was consumed---the fills a
real passive order does not get, and which would have been profitable. The naive backtest is
therefore biased toward adverse-selection-free fills that do not exist.

\subsubsection{Requirements of a queue-exact replay}

A replay that can answer the execution questions of this part needs:
\begin{enumerate}[leftmargin=1.5em]
\item market-by-order data, so that queue position is defined at all;
\item an order-level reconstructed book in which the simulated order occupies a place in the
  first-in-first-out queue;
\item explicit, independent latencies for outbound orders, outbound cancels and inbound market
  data;
\item an explicit fill rule stated in terms of recorded events;
\item a stated assumption about market impact.
\end{enumerate}

\begin{lstlisting}
replay:  strategy sees the book at market time t
         (the feed reached it a feed delay after the exchange event)
         -> order arrives at the exchange at t + feed delay + order delay
         -> appended at the tail of the real FIFO queue at its price
         -> filled only when recorded executions reach an order behind it
         -> P&L from the actual (simulated) fill time and price
\end{lstlisting}

The fill rule rests on price--time priority: if a recorded aggressor reached a resting order that
arrived later than the simulated order at the same price, it must have passed through the
simulated order first. The standard simplifying assumption is \emph{no market impact}: the
recorded market is replayed as it happened, as though the simulated order had not been there. Every
cost measured this way is a cost before the market's reaction to the order.

Replay and generative simulation answer different questions. Generative simulators produce new
market trajectories and are judged on how realistic they are; a replay simulator keeps the recorded
market fixed and produces only the counterfactual outcome of the researcher's own orders. GPU
simulators such as JAX-LOB~\citep{frey2023} move the matching engine onto accelerators to speed up
reinforcement-learning training; the concern here is not throughput but the fidelity of the fill
inference. Section~\ref{sec:kaspar} describes one implementation of a queue-exact replay built to the five
requirements above.

\subsubsection{Limits of replay fidelity}
\label{sec:replay-limits}

A queue-exact replay is exact only with respect to what the feed shows. Two gaps remain.

\paragraph{Hidden liquidity.} Iceberg orders make executable liquidity exceed displayed liquidity,
so a model trained on displayed volume misestimates depletion, impact and fill probability. In
E-mini S\&P~500 futures, one detection study estimates that about 9\% of book volume is iceberg
orders~\citep{christensen2013}; when participants detect an iceberg they trade against it actively, and
icebergs create an adverse-selection cost for other limit orders~\citep{frey2017iceberg}. In a replay, an
iceberg's hidden size is invisible until it trades, so the queue ahead of a simulated order can be longer than the
reconstructed one, and fills are overstated by an amount that cannot be measured from displayed data alone.

\paragraph{Feed and timestamp errors.} At high frequency, data quality is part of the model:
timestamp synchronisation, packet loss, sequence gaps, book reconstruction, exchange-specific event
semantics, stale snapshots, duplicates, and clock differences between market-data and execution
systems.

\begin{example}[A sequence gap]
On a feed with sequence numbers, a missed packet leaves the reconstructed book wrong until the
handler recovers from a snapshot. A model trained on rows from that interval learns
reconstruction error, not market behaviour, and a replay that runs through it fills orders against a book that never
existed. No number of parameters compensates for it; the gap
has to be detected and the interval excluded or repaired.
\end{example}

\subsection{Existing simulators compared}
\label{sec:sim-survey}

This section compares existing simulators in two steps. Table~\ref{tab:simulators} covers the simulators most used
in research, of all four kinds described below; Table~\ref{tab:replay-tools} then narrows to historical-replay tools
only, the kind Kaspar belongs to, and includes every such tool we could classify. A few tools, such as hftbacktest
and NautilusTrader, appear in both.

\paragraph{Four kinds of simulator.} Simulators differ first in where the market comes from.
\begin{itemize}[leftmargin=1.5em,itemsep=1pt]
\item \textbf{Historical replay.} The recorded market is played back event by event, and the researcher's orders are
  matched against it. The market is real but cannot react to those orders.
\item \textbf{Agent-based.} Many simulated traders, each following a rule, send orders to a simulated exchange. The
  market reacts, but it is only as realistic as the agents.
\item \textbf{Model-based.} A fitted stochastic model of the book---such as the queue-reactive model of
  Section~\ref{sec:qr}---generates order flow. The market reacts mechanically, through the model's rates.
\item \textbf{Generative.} A network trained on order data generates order flow, often conditioned on the
  researcher's orders (Section~\ref{sec:deep-gen}).
\end{itemize}
\citet{jain2024sim} review the model-based and generative approaches and the stylised facts used to validate them.
For execution questions the type matters less than five properties: whether the simulated order has an exact
position in a queue of individual orders; whether latency is modelled; whether the simulated order changes the
market; whether a run can be repeated exactly; and whether the same code can trade live.

\paragraph{Simulators of all four kinds.} Table~\ref{tab:simulators} compares the simulators most used in research,
as documented in their papers and code.
\begin{itemize}[leftmargin=1.5em,itemsep=1pt]
\item \textbf{ABIDES}~\citep{byrd2020abides} and its reinforcement-learning interface
  \textbf{ABIDES-Gym}~\citep{amrouni2021abidesgym} are discrete-event agent-based simulators in Python. The exchange
  keeps every order in first-in-first-out queues, and messages between agents and the exchange carry modelled
  network and processing delays.
\item \textbf{MAXE}~\citep{belcak2022maxe} is an agent-based simulator in C++ with a Python interface, with
  price--time and pro-rata matching and message and processing delays.
\item \textbf{PyMarketSim}~\citep{mascioli2024pymarketsim} is an agent-based environment in Python for
  reinforcement learning.
\item \textbf{JAX-LOB}~\citep{frey2023} replays Nasdaq order messages (LOBSTER) on a GPU, so that many simulations run
  in parallel for reinforcement-learning training. It starts each episode from a snapshot in which each price
  level is a single synthetic order, so the agent's position relative to orders that were resting at the start is
  approximate, and it re-submits the agent's orders at each step.
\item \textbf{mbt\_gym}~\citep{jerome2023mbtgym} simulates the price and fills from parametric models, without a
  book; fills depend on the distance of the quote from the mid.
\item \textbf{hftbacktest}~\citep{hftbacktest} replays market-by-price or market-by-order data, mainly from
  cryptocurrency venues. With market-by-order data it keeps every order in a first-in-first-out queue and fills
  the simulated order when an execution reaches an order behind it---the same rule as Kaspar; with price-level
  data it estimates the volume ahead. It models feed latency from the data's timestamps and order latency from
  recorded live latencies, and the same Rust strategy code can trade live on two cryptocurrency exchanges.
\item \textbf{NautilusTrader}~\citep{nautilustrader} is an event-driven backtester and live-trading platform in Rust
  and Python. It tracks an estimate of the displayed quantity ahead of a simulated order, models order-entry
  delays, and runs the same strategy code in backtest and live trading.
\item \textbf{MarS}~\citep{li2025mars} generates order flow with a large generative model trained on Chinese
  equity orders and matches it in an order-level exchange; the generated market responds to injected orders.
\item \textbf{TRADES}~\citep{berti2025trades}, released in the DeepMarket framework, generates orders with a
  diffusion model inside ABIDES.
\item \textbf{The queue-reactive model}~\citep{huang2015} and its neural extension~\citep{deepqr} are model-based
  simulators of the queues at the first five levels (Section~\ref{sec:qr}).
\end{itemize}

\begin{table}[!htbp]
\centering\scriptsize
\caption{Order-book simulators compared on the properties execution research needs, from their papers, documentation
and code; citations are in the list above. Queue position: \emph{exact (recorded)} means the simulated order sits in a reconstructed queue of the
recorded individual orders; \emph{exact (simulated)} means a queue of individual orders in a simulated market;
\emph{estimated} means a model of the volume ahead. ``--'': not provided.}
\label{tab:simulators}
\setlength{\tabcolsep}{2.5pt}
\begin{tabularx}{\linewidth}{>{\raggedright\arraybackslash}p{0.13\linewidth}>{\raggedright\arraybackslash}p{0.12\linewidth}>{\raggedright\arraybackslash}X>{\raggedright\arraybackslash}X>{\raggedright\arraybackslash}p{0.1\linewidth}>{\raggedright\arraybackslash}X>{\raggedright\arraybackslash}p{0.11\linewidth}}
\toprule
Simulator & Market from & Own-order queue position & Latency model & Market reacts & Repeatable runs & Same code live \\
\midrule
ABIDES & agents & exact (simulated) & network and processing delays & yes & same seed, same run (ties within a nanosecond in arbitrary order) & -- \\
ABIDES-Gym & agents & exact (simulated) & pairwise delays & yes & seeded & -- \\
MAXE & agents & exact (simulated); price--time or pro-rata & message and processing delays & yes & not seeded & -- \\
PyMarketSim & agents & price--time book; agents withdraw orders on arrival & -- & yes & partly seeded & -- \\
JAX-LOB & replay (Nasdaq) & approximate: initial levels are one synthetic order each; agent orders re-queued each step & one observation-to-action delay (default 0) & -- & random keys; no claim & -- \\
mbt\_gym & model & no book; fill probability from quote depth & -- & parametric impact & seeded & -- \\
hftbacktest & replay (mainly crypto; CME via vendor files) & exact (recorded) with order-level data; estimated with price-level data & feed latency from data; order latency constant or from recorded live latencies & -- & no claim found & yes (two crypto venues) \\
NautilusTrader & replay & estimated volume ahead & order-entry delays & -- & deterministic runtime claimed; seeded fill model & yes \\
MarS & generative model & exact (simulated) & computation and communication delays & yes & seeded sampling & -- \\
TRADES & generative model in ABIDES & exact (simulated) & -- (as shipped) & yes & seeded & -- \\
Queue-reactive & model & one tagged order, under stated assumptions & -- & yes, through queue sizes & -- & -- \\
\midrule
Kaspar (Section~\ref{sec:kaspar}) & replay (any venue with a feed handler; CME MDP~3.0 implemented) & exact (recorded) & separate order, cancel and feed delays & -- & single thread, market-time clock, seeded; argued, test not reported & live path shares all code up to one branch \\
\bottomrule
\end{tabularx}
\end{table}

\paragraph{What the comparison shows.} Agent-based, model-based and generative simulators let the market react to
the simulated orders, which replay cannot; but in them the queue is exact only inside a simulated market, so fill
rates are only as realistic as the simulated order flow. Among replay tools, exact queue position in the recorded
market requires order-level data and an order-level book: in Table~\ref{tab:simulators}, hftbacktest with
order-level data and Kaspar provide it, JAX-LOB approximates it, and NautilusTrader estimates it. Only hftbacktest,
NautilusTrader and Kaspar run the same strategy code in simulation and live trading; the venues each of them connects
to are compared at the end of this section.

\paragraph{Historical-replay tools only.} Kaspar is a historical-replay tool, so Table~\ref{tab:replay-tools}
compares it with the other tools that replay recorded order-book data and fill a researcher's orders against it. It
lists every such tool we found with enough documentation or code to classify it, including several not in
Table~\ref{tab:simulators}. The list below groups them by how they place the simulated order in the queue.
\begin{itemize}[leftmargin=1.5em,itemsep=1pt]
\item \textbf{Exact position in the recorded queue, recorded market unchanged.} Besides hftbacktest and Kaspar,
  \textbf{lobsim}~\citep{lobsim}, a C++20 library, replays order-level data and counts exactly how much recorded
  quantity is ahead of each simulated order; it has optional order latency, states that it is deterministic, and has
  no live-trading path. \citet{moallemi2017} use the same rule as a method on Nasdaq ITCH data, without software: a
  simulated order counts as filled once a later order at its price is filled.
\item \textbf{Exact position, but the simulated orders change the recorded book.} Two open-source engines insert the
  researcher's orders into the replayed order-level book and match them there, so the orders consume liquidity that
  the recorded market then lacks~\citep{chasemetoyer_bte,koteyevlev_mm}. Neither models latency.
\item \textbf{Estimated position.} Most replay tools work from price-level data and estimate the volume ahead:
  the market-making simulators of \citet{spooner2018} and \citet{sadighian2019}, one open-source backtester for
  cryptocurrency data~\citep{evgerher_hft}, the CME execution model of \citet{dixon2018exec}---which notes that
  order-level data would make the position exact---and \citet{vyetrenko2019}, who add latency taken from historical
  profiles and a synthetic market response. A note by \citet{rigtorp2013queue} describes one such
  estimation rule. Among commercial tools, TT's simulation environment states that queue position is estimated
  for all markets, counting trades ahead of the order but not cancellations~\citep{tt_piq,tt_sim_env}; Deltix
  documents order-level simulators whose fills remove liquidity from the replayed book, but not how they place the
  order in the queue~\citep{deltix_exsim,deltix_qo_sheet}.
\item \textbf{No queue at all.} Some tools fill an order only when the price crosses it~\citep{tribeca}. Retail
  platforms replay the depth of book for display but fill without a queue: Sierra Chart fills a limit order when the
  opposite best quote reaches its price, and states that its queue tracking does not apply in
  replay~\citep{sierrachart_replay}; ATAS fills replayed orders at once and in full~\citep{atas_replay}. NinjaTrader,
  Bookmap and MultiCharts replay depth with simulated trading but do not document how replayed orders are
  filled~\citep{ninjatrader_playback,bookmap_replay,multicharts_sim}.
\end{itemize}
Bar-level backtesters, which see only open, high, low and close prices over fixed intervals or the best quotes, are
outside the scope of this comparison: for example LEAN~\citep{lean}, Jesse~\citep{jesse} and the backtesters of
Hummingbot~\citep{hummingbot} and Barter~\citep{barter}.

\begin{table}[!htbp]
\centering\scriptsize
\caption{Historical-replay backtesters at order-book level, from their papers, documentation and code. ``Recorded
market unchanged'': the simulated orders do not alter the replayed book (no market impact). ``--'': not provided;
``not stated'': documentation silent. Commercial entries (below the line) are from vendor documentation; the Deltix
knowledge base requires a login and was not read, and no readable documentation of a replay fill simulator was found
for CQG.}
\label{tab:replay-tools}
\setlength{\tabcolsep}{2.5pt}
\begin{tabularx}{\linewidth}{>{\raggedright\arraybackslash}p{0.15\linewidth}>{\raggedright\arraybackslash}p{0.12\linewidth}>{\raggedright\arraybackslash}X>{\raggedright\arraybackslash}X>{\raggedright\arraybackslash}p{0.1\linewidth}>{\raggedright\arraybackslash}p{0.11\linewidth}>{\raggedright\arraybackslash}p{0.14\linewidth}}
\toprule
Tool & Data & Own-order queue position & Latency & Recorded market unchanged & Same code live & Feeds \\
\midrule
Kaspar & order-level & exact, in the recorded queue & order, cancel and feed delays & yes & yes & CME MDP~3.0 (own handler) \\
hftbacktest & order- or price-level & exact with order-level data; estimated otherwise & feed from data; order constant or recorded & yes & yes (crypto) & crypto; CME via vendor files \\
lobsim & order-level & exact count of recorded quantity ahead & optional order delay & yes & -- & Coinbase, LOBSTER files \\
NautilusTrader & order- or price-level & estimated & order-entry delays & yes & yes & several venues; CME via vendor files \\
Moallemi--Yuan method & order-level (ITCH) & exact (fill when a later order fills) & -- & yes & -- & Nasdaq \\
Backtesting-Engine & order-level & exact; own orders matched in the replayed book & -- & no & not stated & crypto (CoinAPI) \\
Market\-Making\-ML\-Algo & order log & exact; recorded trades re-matched & -- & no & -- & one equity order log \\
\citet{spooner2018} & price-level, 5 levels & estimated: back of queue, proportional cancels & coded but unused & yes & -- & European equities \\
\citet{sadighian2019} & price-level snapshots & estimated: volume ahead less trades & -- & yes & -- & crypto \\
\citet{dixon2018exec} & price-level & estimated (exact with order-level data) & zero in backtest & not stated & -- & CME ES \\
\citet{vyetrenko2019} & price-level & estimated: back of queue & historical latency profiles & no: synthetic response & -- & futures \\
hft-backtesting & price-level & estimated, probabilistic & single delay & yes & -- & BitMEX \\
\midrule
TT Backtesting & book and trades; price- or order-level not stated & estimated: quantity ahead reduced by trades, not by cancels & not stated & not stated & yes (TT algos) & TT-connected exchanges \\
Deltix QuantOffice & price- or order-level & not stated & configurable & no: fills remove replayed liquidity & yes & several asset classes \\
NinjaTrader Playback & price-level; bid and ask sizes set to 1 & not stated & -- & not stated & yes & data-provider dependent \\
Bookmap replay & depth and best quotes & estimate displayed; use in fills not stated & not stated & not stated & yes & recorded feeds; CME futures \\
MultiCharts & levels 1 and 2 & not stated & not stated & not stated & yes & data-provider dependent \\
Sierra Chart replay & trades with recorded best quotes & none: fills when the opposite quote reaches the price & not stated & not stated & yes & CME and other exchanges \\
ATAS replay & trades and order book & none: filled at once, in full & -- & not stated & manual trading only & subscription dependent \\
\bottomrule
\end{tabularx}
\end{table}

\paragraph{Where Kaspar sits.} Of the replay tools found, three place the simulated order in the recorded queue
without changing the recorded market: hftbacktest, lobsim and Kaspar. Of these, lobsim has no live-trading path, and
hftbacktest trades live on two cryptocurrency exchanges. NautilusTrader, which estimates queue position, connects to
several venues and data vendors. Both reach CME data only through a data vendor's files. Kaspar is the only one we found that decodes an exchange's own market-by-order feed (CME MDP~3.0, from live
multicast or packet captures), models order, cancel and feed delays separately, and runs the same code in replay,
paper and live trading against that exchange. We found no open-source or documented commercial backtester that
decodes MDP~3.0 itself. None of the commercial documentation we could read describes filling a simulated order by its
exact place in a recorded order-level queue: the tools estimate the position, fill without a queue, or do not say.

\paragraph{Running many days at once.} A Kaspar replay processes one trading day per run, from one packet capture or
one decoded file, on one thread. Days are independent, so many days can be simulated at the same time as separate
processes on one large machine---for example a large cloud instance---with one process per day and no licence limit on
the number of runs. The other tools differ on this point (Table~\ref{tab:parallel}). The open-source replay tools allow
the same: hftbacktest's examples run backtests in parallel processes~\citep{hftbacktest}, NautilusTrader recommends one
process per backtest~\citep{nautilustrader}, and ABIDES ships a script that starts one replay process per
date~\citep{abides_repo}; JAX-LOB instead runs up to a thousand environments in parallel on one GPU~\citep{frey2023}.
Among the commercial tools, TT's backtests run only on TT's servers, are charged by run time, and cover at most one
trading day each~\citep{tt_backtest_run,tt_backtest_getstarted}; Bookmap and ATAS allow one running copy or session per
licence or account~\citep{bookmap_sysreq,atas_auth}; and a MultiCharts licence runs online on one computer at a
time~\citep{multicharts_license2pc}. Sierra Chart is the exception: its licence allows many copies on one computer at the
same time~\citep{sierrachart_purchase}. For NinjaTrader's depth replay, Deltix and lobsim the documentation does not
say. Compared with the commercial replay tools, therefore, the ability to run many days in parallel on one's own
hardware is an advantage of Kaspar shared only by Sierra Chart; among open-source tools it is common.

\begin{table}[!htbp]
\centering\footnotesize
\caption{Running many replay simulations at once on one's own hardware, from documentation, licence terms and code.}
\label{tab:parallel}
\begin{tabularx}{\linewidth}{>{\raggedright\arraybackslash}p{0.2\linewidth}>{\raggedright\arraybackslash}p{0.2\linewidth}>{\raggedright\arraybackslash}X}
\toprule
Tool & Parallel runs on own hardware & Restriction or evidence \\
\midrule
Kaspar & yes & one process per day; no licence limit \\
hftbacktest & yes & examples use parallel processes \\
NautilusTrader & yes & one backtest per process \\
ABIDES & yes & script starts one replay process per date \\
JAX-LOB & on one GPU & up to 1,000 parallel environments; GPU memory \\
lobsim & not documented & open licence; no restriction found \\
\midrule
TT Backtesting & no & runs on TT's servers only; one trading day per backtest; charged by run time \\
Bookmap & no & one running copy per licence \\
ATAS & no & one active session per account \\
MultiCharts & limited & online on one computer; a second computer offline for backtesting \\
Sierra Chart & yes & many copies on one computer under one licence \\
NinjaTrader (depth replay) & not documented & one replay connection per session described \\
Deltix QuantOffice & not documented & cloud product lists a parallel runner; licence terms behind a login \\
\bottomrule
\end{tabularx}
\end{table}

\subsection{Kaspar HFT}
\label{sec:kaspar}

This section describes Kaspar, the replay system used as the worked example in this survey: its data path, how it
fills simulated orders, which design choices make its replays repeatable, how the same code runs live, a model-free
benchmark strategy, and how it measures latency and execution cost.

Kaspar is a trading and simulation system written in C++20 on an \emph{actor} framework~\citep{maciejewski_actor}:
every component is an actor, which owns its state and communicates with other actors only by messages. Its order
book, simulated order manager, strategy actors and replay are not tied to one exchange: a venue is connected through
a market-data handler, which turns the venue's feed into book updates, and an order-entry gateway. The handlers
implemented today are for CME---the MDP~3.0 market-by-order feed and the iLink~3 order-entry protocol---and are used
for ES, NQ and Treasury futures. Another market---equities, other futures exchanges, cryptocurrency venues---needs
only its own pair of handlers; the rest of the system is reused unchanged. Every statement below describes the design
and code; the few numbers quoted come from previously published measurements, attributed where they appear.

The data path is
\begin{center}\small
exchange multicast (CME MDP~3.0 today) $\to$ socket reader $\to$ message buffer $\to$ decoder $\to$ feed handler
$\to$ order book $\to$ strategy $\to$ simulated order manager $\to$ [exchange gateway $\mid$ simulated fill].
\end{center}
\emph{Multicast} means the exchange sends each packet once and every subscribed host receives it. The strategy
actors are called \emph{lights} in the code and in the figures. Actors can be placed in a \emph{group}, which runs
several actors on one thread, processing one message at a time in arrival order. The simulated order manager (SOM)
sits between the strategy and the venue. In live mode it forwards orders to the exchange gateway (CME iLink~3); in
simulation mode it sends them into the reconstructed book, where they are matched against recorded order flow.

\begin{figure}[ht]
\centering
\begin{tikzpicture}[
  node distance=4mm and 6mm,
  box/.style={draw, rounded corners=2pt, minimum height=7mm, align=center, font=\small},
  arr/.style={-{Stealth[length=2mm]}, thick}]
\node[box] (md) at (0,0.7) {recorded MBO\\(replay)};
\node[box] (live) at (0,-0.7) {live MDP 3.0\\multicast};
\node[box] (ob) at (3.2,0) {order book\\(per order)};
\node[box] (strat) at (5.6,0) {strategy\\(lights)};
\node[box] (som) at (7.6,0) {SOM};
\node[box] (sim) at (10.6,0.9) {fill in\\reconstructed queue};
\node[box] (ilink) at (10.6,-0.9) {iLink 3\\gateway};
\draw[arr] (md) -- (ob);
\draw[arr] (live) -- (ob);
\draw[arr] (ob) -- (strat);
\draw[arr] (strat) -- (som);
\draw[arr] (som) -- node[above left, font=\scriptsize]{sim} (sim);
\draw[arr] (som) -- node[below left, font=\scriptsize]{live} (ilink);
\draw[arr, dashed] (sim.north) |- ++(0,5mm) -| node[pos=0.25, above, font=\scriptsize]{simulated order enters the book} (ob.north);
\end{tikzpicture}
\caption{Kaspar data and order paths. Everything downstream of the market-data source is the same
code in replay and live operation; the branch between simulated and live execution is a single
decision in the simulated order manager.}
\label{fig:kaspar}
\end{figure}


\subsubsection{Queue-exact simulated fills}
\label{sec:kaspar-fills}

The central requirement of Section~\ref{sec:replay} is that a simulated order should be filled only if
a real order in its place would have been. This section explains how Kaspar meets it: simulated
orders sit in the same per-order queue as real ones, reach the book after a modelled delay, and are
filled only by real executions that reach them. Here and throughout, \emph{queue-exact} means exact
with respect to the displayed order sequence carried by the input MBO feed and the documented matching
rules; it does not imply knowledge of hidden liquidity or unobserved participants, and it does not model
the market impact of the simulated orders.

\paragraph{The book holds every order.}

The simulator reconstructs a full market-by-order book from recorded exchange data. Each price
level is a queue that holds every resting order individually, in arrival order, and simulated
orders are held in the same queue alongside real ones. There is no estimated ``volume ahead''
variable: a simulated order's queue position is its position in the list. This is the property
that makes a fill inference definable at all~\citep{maciejewski_shadow}.

\paragraph{Placement and latency.}

In simulation mode the SOM turns a new order into a book \emph{add} message tagged with a
simulated venue identifier and sends it to the order-book actor. The book does not insert it
immediately. It holds the order on a delay queue until its modelled arrival time at the matching
engine:
\begin{lstlisting}[language=C++]
auto order_engine_arrive_time =
    order_leave_time
      + uint64_t(feed_delay) * 1000
      + uint64_t(std::max(40, eff_delay)) * 1000; // 40 us floor
\end{lstlisting}
The outbound order delay, the outbound cancel delay and the inbound feed delay are independent
parameters. The feed delay is added because the strategy saw the market late by that amount, so a
round trip is feed plus order, not order alone. When the order arrives it is appended at the tail
of the queue at its price, behind all depth already resting there.

\paragraph{The fill rule.}

A recorded cancellation of a real order simply unlinks it, and every order behind it moves up. A
recorded execution triggers simulated fills only when the executing real order is the first real
order in the queue and a simulated order sits ahead of it:
\begin{lstlisting}[language=C++]
// fill sim orders only if the executing order is at the front
if (o == q->get_head_no_sim())
{
  fill_prev_sim_order();
  ...
}
\end{lstlisting}
Any remaining executed quantity is carried forward to earlier simulated orders. This is the
price--time-priority argument of Section~\ref{sec:replay} in code: an aggressor that reached a
real order which arrived after the simulated order must have passed through the simulated order.
Simulated orders are excluded from the best bid and offer that the strategy sees, so a strategy
cannot react to its own quote. Market impact is not modelled~\citep{maciejewski_shadow}.

\begin{example}[A fill, step by step]
\label{ex:trace}
Bid queue at 5000.00 over seven successive steps \#0--\#6, sizes in lots, real orders R1--R4 and
one simulated order S:
\begin{lstlisting}
#0  resting:            R1(10) R2(5) R3(20)
#1  S(3) arrives:       R1(10) R2(5) R3(20) S(3)      appended at tail
#2  R2 cancels:         R1(10) R3(20) S(3)            S moves up
#3  exec 10 vs R1:      R3(20) S(3)                   none ahead of R1: no fill
#4  exec 20 vs R3:      S(3)                          S is behind R3: no fill
#5  R4(8) joins:        S(3) R4(8)
#6  exec 5 vs R4:       R4 is head real, S ahead  ->  S filled 3 @ 5000.00
\end{lstlisting}
At step \#6 the aggressor reached R4, which arrived after S, so it must have passed through S. A
``touch'' backtest would have filled S at step \#3, three recorded events earlier and before the
queue ahead of it was consumed.
\end{example}

\paragraph{Market impact: a limitation of replay, and a possible extension.} Because the recorded
market behaves exactly as it did, whatever the simulated orders do, replay answers ``would my order have
been filled against this market?'' but not ``how would the market have moved because of my order?''. A
simulated order never consumes liquidity that real traders would have reacted to. Market impact is
therefore outside what this architecture measures, and it is not implemented.

One way to add a reacting market, which we describe as possible future work rather than a feature, is a
second market-data source: a queue-reactive simulator (Section~\ref{sec:qr}) in place of the recorded feed.
The separation between the source of market data and the book, order manager and strategy actors means
the rest of the pipeline would not change. The generator produces queue counts rather than individual
orders, so its events would have to be turned into synthetic order-level messages---order identifiers and
first-in-first-out placement---for the existing book. The researcher's orders would then change queue
sizes, the simulator's rates would respond, and the price response to the strategy's own trading could be
measured.

What such an extension would estimate is limited, and should be stated with it.
\begin{itemize}[leftmargin=1.5em,itemsep=1pt]
\item It would capture \emph{mechanical} impact: queues consumed by the strategy's orders, refilling at
  state-dependent rates, and price moves when a queue empties. Neural extensions of the queue-reactive
  model reproduce the square-root law of impact on Bund futures~\citep{deepqr}, so this part can be
  realistic.
\item It would not capture \emph{informational} impact: other participants inferring that a large buyer is
  present and repricing. In the queue-reactive model this sits in the random resets of the book, which do
  not respond to the strategy.
\item Its impact would be a property of the simulator, and would need checking against impact measured on
  real trading---post-fill mark-outs from paper or live fills, or published impact laws---before being
  relied on.
\end{itemize}
The resulting division of labour would be: queue-exact replay for fill realism on the real market, without
impact; a reactive simulator for mechanical impact and for comparing tactics in a market that responds;
and paper and live trading as the ground truth for both.


\subsubsection{Deterministic replay}
\label{sec:kaspar-determinism}

In the replay configuration the file reader, the order book, the timer, the SOM, the strategy
actors and the position manager all run in one actor group on one thread. The reader is added to
the group last, so all consumers exist before the first market-data message is delivered. Three
design choices then make a replay repeatable:
\begin{enumerate}[leftmargin=1.5em]
\item \textbf{One message sequence.} A single thread processes every message---market data, orders,
  acknowledgements, fills, timer events---in one total order.
\item \textbf{Market time.} The timer actor is driven by market-data timestamps, not the wall
  clock, so timeouts fire at the same point in the data on every run.
\item \textbf{Seeded randomness.} Any randomised strategy decision uses a per-actor seeded
  generator set in configuration.
\end{enumerate}
The property this implies, and the test that would check it, is that replaying the same recorded
file twice with the same configuration yields byte-identical fill records. We argue the property
from the design; we do not report the test here.

Running the same program in research and in production also removes one source of difference between them: a
result cannot change because the production code was written separately.


\subsubsection{Replay, paper and live trading from the same code}
\label{sec:kaspar-equivalence}

Two choices are independent: where market data come from (recorded file or live multicast), and
where orders go (reconstructed book or exchange). Everything downstream of the message buffer---the
decoder, feed handler, order book, strategy and SOM---is the same code whichever source feeds it.
The branch between simulated and live execution is a single condition in the SOM's new-order
handler; both paths apply the same pre-trade reject checks, with one additional rule in live mode
for cancels of orders not yet acknowledged by the exchange. The combination ``live data, simulated
execution'' is paper trading: orders enter the live reconstructed book and are filled by the live
order flow under the same rule as in replay.

The live order-entry path implements CME's iLink~3 protocol: binary message encoding, HMAC request
signing (a keyed checksum that authenticates each message), sequence management, and failover between primary and secondary gateways. The project
documentation states that both the market-data handler and the order-entry session have passed CME's
automated certification. Strategies can be written as in-process actors in C++ or, through an
in-process foreign-function interface, in Rust.

Two qualifications keep the claim precise. Switching between simulated and live execution is a single
flag in the SOM, but in the distributed build that flag is set in the program source rather than in a
configuration file, and the distributed configuration runs in simulation mode; the equivalence
demonstrated day to day is therefore between replay and paper trading, with the live path sharing all
code up to the SOM branch. And equivalence of code is not equivalence of outcomes: replay assumes no
market impact (Section~\ref{sec:kaspar-fills}) and models latency with configured delays, so a strategy
that is profitable in replay can still behave differently live.


\subsubsection{A model-free execution benchmark}
\label{sec:shadow}

\pov{} is the passive execution algorithm used in the case study~\citep{maciejewski_shadow}. It
makes no prediction about book evolution. When a third-party limit order is added within a
configured distance of the touch, the algorithm may place its own order at the same price, records
that participant's exchange order identifier, and cancels when that order leaves the book. The
placement decision is randomised at a configured rate; cancellation can be delayed by a grace
period measured in events or time. The design comment for the grace period states the trade-off
of Section~\ref{sec:adverse} directly: cancel too quickly and fills are given up; cancel too
slowly and the order bears the adverse selection.

Its role in this survey is as a null model for passive placement. It defines the bottom rung of a
hierarchy:
\begin{center}
model-free placement $\;\to\;$ predictive placement $\;\to\;$ full predictive execution policy.
\end{center}
A predictive model that claims value for passive execution should beat the model-free benchmark
\emph{in the same replay, under the same fill rule}, on slippage and on post-fill mark-outs. Any
improvement measured against a naive backtest instead is confounded by the fill-inference bias
of Section~\ref{sec:replay}.


\subsubsection{Latency and slippage measurement}
\label{sec:kaspar-measurement}

Latency claims need their endpoints. In the published measurement of Kaspar's market-data
handler, the start point is a software timestamp at the socket read and the end point is the order
book being published; time spent in the network card and kernel before the read is outside the
measurement, which is therefore socket-to-book, not wire-to-book. Measured on live CME multicast,
the median is 8.0\,$\mu$s, with book-message medians of 7.1, 7.3 and 7.1\,$\mu$s for ES, NQ and ZN
and 99th percentiles of 18.5, 13.6 and 57.0\,$\mu$s~\citep{maciejewski_mdlatency}. These numbers are
quoted as an example of stating endpoints, not as a result of this paper.

The simulator's slippage instrumentation measures execution cost (\emph{slippage}) for each fill against four
references: the mid at decision time, the interval volume-weighted average price, same-side
passive fills by other participants, and the touch. Comparing against the interval VWAP separates
market drift from adverse selection. Post-fill mark-outs at fixed times and after a fixed number
of book events provide the mark-out $R^{\text{fill}}_{h}$ of~\eqref{eq:markout} (the event-count version is
its analogue on an event clock) and the adverse-selection quantity of Section~\ref{sec:adverse}.

\subsection{Adding an order-book model to Kaspar}
\label{sec:kaspar-actor}

Any model of Part~II can be added to Kaspar as one more actor between the order book and the strategy. The actor
subscribes to the book, computes its prediction each time the book changes, and sends the prediction to the
strategy actors as a message (Figure~\ref{fig:kaspar-actor}). Because it is an actor like the others, the same model
runs in replay, paper trading and live trading. The description below follows the code; the parts marked
\emph{new} do not exist in the repository and would have to be written.

\begin{figure}[!htbp]
\centering
\begin{tikzpicture}[font=\scriptsize, scale=0.86, transform shape,
  a/.style={draw, rounded corners=2pt, align=center, minimum height=10mm, text width=22mm},
  n/.style={a, fill=orange!15, very thick, text width=28mm},
  arr/.style={-{Stealth[length=1.6mm]}, thick}]
\node[a, fill=gray!10] (feed) at (0,0) {market data\\(replay or live)};
\node[a, fill=blue!10] (ob) at (3.3,0) {order book\\actor (OB)};
\node[n] (pred) at (8.8,0) {\textbf{model actor} (new):\\queue imbalance, a Hawkes\\intensity, a network};
\node[a, fill=green!12] (light) at (13.4,0) {strategy actors\\(lights)};
\node[a, fill=gray!10] (som) at (13.4,-2.2) {simulated order\\manager (SOM)};
\draw[arr] (feed) -- (ob);
\draw[arr] (ob) -- node[above]{\texttt{EndOfBurst}} node[below, align=center]{book, 16 levels,\\after each update} (pred);
\draw[arr] (pred) -- node[above]{\texttt{Prediction}} node[below]{(new message)} (light);
\draw[arr] (ob.north) to[out=35, in=145] node[above]{\texttt{EndOfBurst} (existing subscription)} (light.north);
\draw[arr] (light) -- node[right]{orders} (som);
\draw[arr] (pred.south) -- ++(0,-0.9) -| node[pos=0.25, below]{\texttt{Subscribe(HI)}, sent once at start} (ob.south);
\end{tikzpicture}
\caption{An order-book model added to Kaspar as an actor. It subscribes to the order book, receives the book after
each burst of updates, and sends a prediction message to each strategy actor. Message names are those in the code;
the model actor, the \texttt{Prediction} message and its handler in the strategy are new.}
\label{fig:kaspar-actor}
\end{figure}

\begin{enumerate}[leftmargin=1.8em,itemsep=2pt]
\item \textbf{Define the message.} A new message type derives from \texttt{actors::MessageT<Prediction>}, which
  assigns it a unique identifier automatically; it carries the instrument, the market time of the book update it
  was computed from, and the predicted value.
\item \textbf{Define the actor.} The model actor derives from \texttt{actors::Actor} and registers two handlers
  with \texttt{MESSAGE\_HANDLER}: one for \texttt{Start}, sent by the framework when the system starts, and one for
  \texttt{EndOfBurst}, the message the order book sends after each burst of updates. In its \texttt{Start} handler it
  sends \texttt{Subscribe(HI)} to the order book---the same call the strategy actors and the simulator's slippage
  probe already make.
\item \textbf{Compute and publish.} On each \texttt{EndOfBurst} it reads the prices and sizes of up to 16 levels
  from the message, computes the prediction, and sends a \emph{new} \texttt{Prediction} message to each strategy
  actor; the framework does not allow one message object to be sent twice.
\item \textbf{Use the prediction in the strategy} (new). The strategy actor registers a handler for
  \texttt{Prediction}, stores the latest value, and uses it in its placement decision---for example as an extra
  condition before it places an order.
\item \textbf{Wire it in.} In the replay simulator the actor is added to the single actor group before the
  market-data reader, so that it exists before the first message arrives; in the trading binary it is added
  alongside the strategy actors of the same order book.
\end{enumerate}

In outline (the full sketch, with the file and line of every call used, is in the notes accompanying this survey):
\begin{lstlisting}[language=C++]
namespace ob_msg  = frame::ob::msg;
namespace mda_msg = frame::mda::msg;

// New message type: the model's output for one book update.
struct Prediction : actors::MessageT<Prediction> {
  uint     sym;          // instrument
  uint64_t market_time;  // time of the book update it came from
  double   value;        // e.g. queue imbalance in [-1, 1]
  Prediction(uint s, uint64_t t, double v)
      : sym(s), market_time(t), value(v) {}
};

// New actor: subscribes to the book, publishes one prediction per update.
class ImbalanceModel : public actors::Actor {
  actor_ptr ob;                    // the order book it reads
  std::vector<actor_ptr> lights;   // the strategy actors it feeds

 public:
  ImbalanceModel(actor_ptr ob_, std::vector<actor_ptr> lights_)
      : ob(ob_), lights(std::move(lights_)) {
    MESSAGE_HANDLER(actors::msg::Start, on_start);
    MESSAGE_HANDLER(ob_msg::EndOfBurst, on_book);
  }

  void on_start(const actors::msg::Start*) {
    ob->send(new mda_msg::Subscribe(mda_msg::Subscribe::HI), this);
  }

  void on_book(const ob_msg::EndOfBurst* m) {
    const auto& pl = m->payload;          // the book after the burst
    // skip bad data and bursts sent while the book is being rebuilt
    // after a gap, as the strategy actors do
    if (!pl || pl->point_.baddata || pl->recovery) return;
    double b = pl->point_.bid_sz[0];      // size at the best bid
    double a = pl->point_.ask_sz[0];      // size at the best ask
    if (a + b <= 0) return;
    double imb = (b - a) / (b + a);
    for (auto l : lights)                 // a new message per recipient
      l->send(new Prediction(pl->sym, pl->txtim_epoch, imb), this);
  }
};
\end{lstlisting}

Three properties of the present code matter for anyone doing this.
\begin{itemize}[leftmargin=1.5em,itemsep=1pt]
\item \textbf{Ordering and determinism.} In the replay simulator every actor runs in one group, on one thread, in
  one message order, so a replay with the model actor added is still deterministic. From reading the code (not
  tested), a prediction sent with \texttt{send} for a book update reaches the strategy after the strategy has acted
  on that update and before the next one; the framework's synchronous \texttt{fast\_send} would deliver it first.
  In the trading binary each actor has its own thread, and the order in which the book update and the prediction
  reach the strategy is not fixed.
\item \textbf{Which book.} The model actor should subscribe to the price-level book (OB). The order-level book
  (TachBook) accepts high-priority subscribers only under specific names, and in the trading binary it receives
  data only in a measurement mode that starves the other consumers.
\item \textbf{What is missing.} There is no general publish--subscribe bus: a model actor must hold a pointer to
  the book and to each strategy actor. The strategy has no input for an external signal today; its only numeric
  input sets a position target.
\end{itemize}

% ================================================================
\section{Hardware acceleration of order-book models}
\label{sec:hardware}
% ================================================================

Hardware bears directly on the second organising principle, Occam's razor. The economic value of a signal
depends on the latency and the determinism of the infrastructure that evaluates it---how long evaluation takes, and
how much that time varies from one event to the next (\emph{jitter})---and richer models tend to raise both latency
and jitter:
\begin{equation}
  \text{model complexity} \uparrow \;\Rightarrow\; \text{inference latency} \uparrow \;\Rightarrow\;
  \text{signal value at execution} \downarrow .
\end{equation}
\inwords{a bigger model takes longer to evaluate, and by the time it answers, some of what it
predicted has already happened.}

This section first follows one market event through a model-driven trading system and shows where the time is
spent; it then explains why recurrent networks and Transformers are slow at that task, collects the published and
the independently audited latency measurements, and lists the stages worth accelerating. The remaining subsections
review the hardware---field-programmable gate arrays (FPGAs: chips whose circuits are configured after manufacture to
implement one design), graphics processors (GPUs), accelerators between the two, and the network---and close with a
guide to placing each workload and a limit on what hardware can achieve.

\subsection{Stages of the per-event pipeline}
\label{sec:hw-pipeline}

This section follows one market event through the system, separates the stages fixed by the exchange protocol from
the stage that depends on the model, and summarises what is known about the cost of each.

Between a market-data packet arriving and an order leaving, a model-driven system runs a pipeline
(Figure~\ref{fig:hw-pipeline}): decode the exchange message; update the book; compute the model's inputs; evaluate the
model; decide; send the order. Each stage has a different cost structure, and each is accelerated in a different way.

\begin{figure}[!htbp]
\centering
\begin{tikzpicture}[font=\scriptsize, scale=0.92, transform shape,
  s/.style={draw, rounded corners=2pt, align=center, text width=19mm, minimum height=13mm},
  arr/.style={-{Stealth[length=1.5mm]}, thick}]
\node[s, fill=blue!10] (a) at (0,0) {\textbf{decode}\\packet $\to$ message fields};
\node[s, fill=blue!10] (b) at (2.4,0) {\textbf{book update}\\find level and order, change it};
\node[s, fill=green!12] (c) at (4.8,0) {\textbf{features}\\window of snapshots, OFI, \ldots};
\node[s, fill=orange!20] (d) at (7.2,0) {\textbf{model}\\forward pass};
\node[s, fill=gray!10] (e) at (9.6,0) {\textbf{decide}\\place, keep, cancel};
\node[s, fill=gray!10] (f) at (12.0,0) {\textbf{order out}\\encode, send};
\foreach \x/\y in {a/b,b/c,c/d,d/e,e/f} \draw[arr] (\x) -- (\y);
\node[align=center, text width=40mm, anchor=north] at (1.2,-0.9) {regular, protocol-bound;\\one message at a time};
\node[align=center, text width=22mm, anchor=north] at (4.8,-0.9) {few operations;\\memory access};
\node[align=center, text width=22mm, anchor=north] at (7.2,-0.9) {tens to millions of\\operations, by model};
\node[align=center, text width=40mm, anchor=north] at (10.8,-0.9) {simple logic;\\must be deterministic};
\end{tikzpicture}
\caption{The per-event pipeline of a model-driven trading system. The model stage is the only one whose cost depends
on the modeller's choices. Schematic.}
\label{fig:hw-pipeline}
\end{figure}

\paragraph{Book building and inference.} The stages fall into two groups that are often conflated. \emph{Book
building}---decoding market data and maintaining the order book---is fixed by the exchange protocol, highly regular,
and the same for every model. \emph{Inference}---evaluating a model on the book---depends on the model. A latency figure for one group should not be compared directly with a figure for the other.

\paragraph{Decoding and book update.} FPGAs have long been used for book building. Decoding market-data feeds on FPGAs
is well documented~\citep{leber2011,morris2009fpga}: \citet{leber2011} decode Ethernet, IP, UDP and the
FAST protocol in hardware with a fourfold latency reduction over software, and \citet{lockwood2012} report a fixed
end-to-end latency of 1\,\textmu s at 10\,Gb/s line rate for a library of networking and protocol-parsing blocks.
The book update is a search in an ordered set of price levels, followed by a change to one level or one order; with
standard ordered containers each search takes time proportional to the logarithm of the number of levels, and an
industry preprint reports a specialised structure 9 to 15 times faster than the C++ standard library's ordered map on
real market data~\citep{krapivensky2025glass}. We found no peer-reviewed measurement of the cost of book updates on
CPUs. On FPGAs the book itself has been built in hardware: a depth-of-book design for options averages about
253\,ns per update, estimated from clock-cycle counts~\citep{dvorak2014book}; a later design reports 26 to 209\,ns,
3.5 to 37.6 times faster than its software version~\citep{liu2024book}; and a complete trading system written in
high-level synthesis (C++ compiled into FPGA logic)---network stack, parsing, book and strategy---measured an on-board round trip below
870\,ns~\citep{boutros2017build}. For Kaspar's software path, socket-to-book latency is reported in Section~\ref{sec:kaspar-measurement}.

\paragraph{The model.} This is the stage that grows with model complexity. A linear score on a few order-flow
features takes tens of multiply--add operations; DeepLOB has about 60,000 weights (Table~\ref{tab:deeplob-params}),
and each forward pass uses each of them at least once. Placing models on the same FPGA as the book is more recent:
tool flows for fast neural-network inference on FPGAs~\citep{hls4ml} make small, quantised networks feasible with
very low and fixed latency, but model capacity is bounded by the size of the FPGA's programmable logic. The next
subsection explains why larger models are slower than their operation count suggests.

\subsection{Sources of per-event latency in recurrent networks and Transformers}
\label{sec:hw-why}

Training processes many examples at once, and graphics processors are built for that. A trading system needs one
prediction for the event that has just arrived---a \emph{batch} of one---and as soon as possible. This section
explains four effects that make that case slow and then collects the published inference times of order-book
models.

\paragraph{1. Recurrence is sequential.} A recurrent network computes its memory at step $t$ from its memory at step
$t-1$ (Section~\ref{sec:deep-families}), so the 100 steps of a window must be computed one after another; this
``inherently sequential nature precludes parallelization within training examples''~\citep{vaswani2017}. A layer needs
a number of sequential steps that grows with the sequence length $n$, where self-attention needs a constant
number~\citep{vaswani2017}. On an FPGA, a recurrent layer built from one reused block has a latency that grows linearly
with the sequence length; unrolling it over time steps removes that growth but multiplies the hardware used by the
sequence length~\citep{khoda2023rnn}.

\paragraph{2. Attention grows with the square of the window.} Self-attention compares every element of the sequence
with every other, so its cost per layer grows as $n^2 d$, where $n$ is the sequence length and $d$ the size of each
element's vector~\citep{vaswani2017}; ``the time and memory complexity of self-attention are quadratic in sequence
length''~\citep{dao2022flashattention}. Doubling the window from 100 to 200 events doubles the work of most layers but
quadruples that of attention.

\paragraph{3. Small batches are limited by memory, not arithmetic.} A processor can only compute as fast as it is fed
with data. The \emph{roofline} model makes this precise: attainable speed is the smaller of the peak arithmetic rate
and the memory bandwidth times the number of operations performed per byte loaded~\citep{williams2009roofline}. With a
batch of one, each weight loaded from memory is used for about one multiplication and one addition, so the
computation is limited by memory bandwidth~\citep{kim2023fullstack}; for large language models ``at small batch sizes
\ldots\ the time to load weights dominates''~\citep{pope2023scaling}. Attention is itself mostly limited by memory
traffic, which FlashAttention reduces by keeping blocks of the computation in on-chip memory~\citep{dao2022flashattention}.
These measurements are for large models; whether weight loading dominates for a model as small as DeepLOB, whose
weights fit in on-chip memory, has not been measured. One measured case shows the batch effect directly: an LSTM on a
V100 GPU processed 660 events per second at batch one, about 7,700 at batch ten and about 30,000 at batch
100~\citep{khoda2023rnn}.

\paragraph{4. Fixed overheads per operation.} Each operation sent to a GPU runs as a separate program, called a
\emph{kernel}, and each kernel launch costs an overhead of the order of microseconds. In NVIDIA's measurement a kernel that ran for 2.9\,\textmu s cost 9.6\,\textmu s per launch when each
launch waited for the previous one, 3.8\,\textmu s when launches overlapped, and 3.4\,\textmu s when the sequence was
recorded once as a \emph{CUDA graph} (CUDA is NVIDIA's GPU programming platform) and replayed~\citep{nvidia_cudagraphs_blog}; graphs pay off only when the same
sequence is reused many times~\citep{nvidia_cuda_guide}. A network of many small layers is many small kernels. Moving
the input to the GPU and the result back crosses the PCIe bus, the link between the processor and add-in cards such
as GPUs and network cards, whose bandwidth is far below the GPU's own memory
bandwidth~\citep{nvidia_cuda_bpg}; a single small transfer over PCIe has been measured at about half a microsecond
(for a network card, not a GPU)~\citep{neugebauer2018pcie}.

\paragraph{Measured latencies of order-book models.} Table~\ref{tab:lob-latency} collects the published inference
times we found. They range from about 25\,\textmu s to about 23\,ms---from about 3 to about 3,000 times the roughly
8\,\textmu s socket-to-book latency of Section~\ref{sec:kaspar-measurement}; about half are below one millisecond
and half above. They are not comparable across rows: devices and software differ, and some
sources do not state the batch size. Latency also does not follow operation counts: in one study a gradient-boosted model with 23
times more structural work than another was 15 times faster~\citep{hedges2026frontier}.

\begin{table}[!htbp]
\centering\footnotesize
\caption{Published inference times of order-book models. Each row is as reported by the source.}
\label{tab:lob-latency}
\begin{tabularx}{\linewidth}{>{\raggedright\arraybackslash}p{0.2\linewidth}>{\raggedright\arraybackslash}X>{\raggedright\arraybackslash}p{0.33\linewidth}}
\toprule
Source & Models and times & Conditions \\
\midrule
\citet{zhang2019} & DeepLOB 0.253\,ms & batch size and device for the timing not stated \\
\citet{tran2019tabl} & C(TABL) 0.025, CNN 0.061, LSTM 0.229\,ms & one sample, Intel i7-4790 CPU \\
\citet{liu2026microsecond} & 8-bit integer 1D-CNN on FI-2010: 57\,\textmu s & FPGA (KU040), inference core only; from the abstract \\
\citet{berti2025tlob} & MLP 0.08, LSTM 0.21, DeepLOB 1.31, TLOB 2.24, MLPLOB 4.79\,ms & RTX 3090 GPU; batch size not stated \\
\citet{hedges2026frontier} & median at batch one: MLPLOB 2.1\,ms (CPU), 2.8\,ms (GPU); TLOB 22.8\,ms (CPU), 5.1\,ms (GPU) & Apple M3 Max, PyTorch, 32-bit floating point; forward pass only, no preprocessing or transfer \\
\bottomrule
\end{tabularx}
\end{table}

\subsection{Audited benchmarks}
\label{sec:hw-audited}

The times in Table~\ref{tab:lob-latency} come from research papers, each under its own conditions. Most published
latency figures for trading hardware come from vendors. The exception is the benchmarks of STAC (the Securities
Technology Analysis Center), run by an independent organisation to a published specification and audited, of which
three bear on this survey (Table~\ref{tab:stac}). Their public summaries are chosen by the vendor that submitted the system, and the full
reports require registration, so the figures below are those summaries.
\begin{itemize}[leftmargin=1.5em,itemsep=1pt]
\item \textbf{STAC-T0} measures network input and output only: the time from the last bit of the market-data packet
  needed for a decision to the first bit of the order, with trivial decision logic.
\item \textbf{STAC-T1} measures a whole \emph{tick-to-trade} path---from a market-data packet arriving to the
  resulting order leaving---on CME E-mini data: decode the feed, keep the state, decide,
  and send the order.
\item \textbf{STAC-ML Markets (Inference)} measures model inference alone, for three sizes of LSTM---the smallest
  under 200,000 parameters, according to one vendor---and later for gradient-boosted trees. It excludes reading the
  data and computing the features. One of its two suites sends each window as a whole; the other, \emph{Tacana},
  slides the window by one step and allows earlier computation to be reused, which is the incremental case discussed
  in Section~\ref{sec:hw-what}~\citep{stac_ml_wg,stac_ml_limits2025}.
\end{itemize}

\begin{table}[!htbp]
\centering\footnotesize
\caption{Audited STAC results relevant to order-book models, from STAC's public summaries. The three benchmarks
measure different paths and must not be compared with one another. ML: 99th-percentile latency, one model
instance, sliding-window (Tacana) suite.}
\label{tab:stac}
\begin{tabularx}{\linewidth}{>{\raggedright\arraybackslash}p{0.13\linewidth}>{\raggedright\arraybackslash}X>{\raggedright\arraybackslash}p{0.37\linewidth}}
\toprule
Benchmark & System & Result \\
\midrule
STAC-T0 & LDA network stack on an FPGA card (2020)~\citep{stac_xlx200514} & minimum 24.2\,ns; mean 27.9--29.2\,ns \\
STAC-T0 & Exablaze ExaNIC V5P (2019)~\citep{stac_exa181106} & 31 to 44\,ns \\
STAC-T1 & ADHOC appliance, AMD Versal FPGA (2024)~\citep{stac_adhc240918} & mean 115\,ns; 99th percentile 140\,ns \\
STAC-T1 & Exegy Xero FPGA card (2019)~\citep{stac_exg191010} & mean 0.552\,\textmu s; 99th percentile 0.609\,\textmu s \\
STAC-ML & NVIDIA A100 GPU (2023)~\citep{stac_ml_nvda2023} & smallest LSTM 35.2\,\textmu s; largest 640\,\textmu s \\
STAC-ML & Myrtle VOLLO, Intel Agilex FPGAs (2023)~\citep{stac_ml_mrtl2023} & smallest LSTM 5.07\,\textmu s; largest 31.0\,\textmu s \\
STAC-ML & NVIDIA GH200 (2025)~\citep{stac_ml_smc2025} & smallest LSTM 4.70\,\textmu s; largest 15.8\,\textmu s \\
STAC-ML & Myrtle VOLLO, AMD Versal FPGA (2026)~\citep{stac_ml_mrtl2026} & smallest LSTM below 2\,\textmu s; largest below 8\,\textmu s \\
STAC-ML & Xelera Silva, AMD Alveo V80 (2026)~\citep{stac_ml_xlra2026} & gradient-boosted trees (different suite): at most 1.95\,\textmu s for the two smaller models, 2.88\,\textmu s for the largest \\
\bottomrule
\end{tabularx}
\end{table}

Three points follow. A complete tick-to-trade path on an FPGA, with simple decision logic, runs in about a tenth to
half a microsecond. Inference of a small LSTM, in the sliding-window setting, now takes single-digit microseconds on
both FPGAs and GPUs: in these results a GH200 GPU was slightly faster than the 2023 FPGA system and slower than the
2026 one. And no audited benchmark yet combines building the book with evaluating a learned model on it, or uses a
published order-book architecture.

\subsection{Stages that need acceleration}
\label{sec:hw-what}

Combining the pipeline of Section~\ref{sec:hw-pipeline}, the four effects of Section~\ref{sec:hw-why} and the
measurements above gives a short list of what to accelerate.
\begin{enumerate}[leftmargin=1.8em,itemsep=2pt]
\item \textbf{Decoding and the book update}, for every model. They are regular and protocol-bound, and FPGAs handle
  them well (Section~\ref{sec:hw-landscape}).
\item \textbf{The input window, incrementally.} A window of the last 100 snapshots changes by one row per event.
  Keeping it as a ring buffer---a fixed array in which the newest row overwrites the oldest---and updating only the
  new row, rather than rebuilding it, removes work that does not
  depend on the model. (Our observation; we found no paper that measures it.)
\item \textbf{The sequential part of recurrent models.} Either shorten the sequence, keep the recurrent state from
  the previous event so that only one new step is computed per event (possible when the window slides by one), or
  move the recurrence to hardware that runs it as a pipeline~\citep{khoda2023rnn}.
\item \textbf{Attention over long windows.} Shorter windows, attention that reuses the keys and values of earlier
  events rather than recomputing them, and memory-aware kernels such as FlashAttention~\citep{dao2022flashattention}.
\item \textbf{Memory traffic and launch overhead at batch one.} Smaller, lower-precision weights: integer
  quantisation (storing weights as low-precision integers instead of floating-point numbers) can cut memory and latency by a factor of 4 to 8 in practice~\citep{gholami2022quant,jacob2018}, and
  pruning plus quantisation compressed a speech LSTM twentyfold with negligible loss of accuracy~\citep{han2017ese}; fusing
  operations into fewer kernels and recording them as a CUDA graph~\citep{nvidia_cudagraphs_blog}; or keeping the whole
  model on the device that receives the data, to avoid PCIe transfers.
\end{enumerate}
We found no published FPGA implementation of DeepLOB, TransLOB or TLOB with a measured latency. The closest are an
8-bit convolutional network for the FI-2010 data at 57\,\textmu s on an FPGA~\citep{liu2026microsecond}, a state-space
model designed for an FPGA from the start, reported at 247\,\textmu s end to end on cryptocurrency order
books~\citep{popov2026ssm}, and the audited LSTM results of Table~\ref{tab:stac}, which reach a few microseconds
for small generic LSTMs.

\subsection{Determinism as a risk-control property}

The figures above are averages and percentiles; the variation of latency matters as well as its typical value.
General-purpose processors, even with \emph{kernel-bypass} networking (packets passed between the network card and
the application without going through the operating system), introduce jitter from caches, branch prediction,
interrupts and scheduling. A strategy that is profitable at its average latency
can lose money if latency spikes at the moments when its signal is acted on.
Deterministic, fixed-cycle execution is therefore a risk-control property, not only a speed
optimisation.

\subsection{Commercial FPGA products for trading}
\label{sec:hw-landscape}

Most production hardware for trading is commercial, and most of its published figures are vendor
claims rather than audited measurements. This section lists representative products and how to read their figures. Table~\ref{tab:vendors} lists representative offerings with
the figures their makers publish; audited results are in Section~\ref{sec:hw-audited}, and the STAC-T0 result of an
Exegy framework on an AMD card is shown here for reference~\citep{exegy_stac2024}. Two cautions apply when reading such
numbers. They measure different things---a transceiver loopback (a signal sent out and straight back through the circuit
that puts bits on the wire), a trigger, a book build, a model
inference---that must not be compared directly. And the AMD transceiver figure explicitly excludes
protocol processing and programmable logic, so it is a lower bound on a component, not a
tick-to-trade time.

\begin{table}[!htbp]
\centering\footnotesize
\caption{Representative commercial hardware for trading and the figures published by the vendors.
Except for the STAC result, figures are vendor claims without independent audit.}
\label{tab:vendors}
\begin{tabularx}{\linewidth}{>{\raggedright\arraybackslash}p{0.2\linewidth}>{\raggedright\arraybackslash}X>{\raggedright\arraybackslash}X}
\toprule
Offering & What it is & Published figure \\
\midrule
AMD Alveo UL3524 / UL3422~\citep{amd_ul3524,amd_ul3422} & FPGA cards built for trading, 64 low-latency transceivers & under 3\,ns transceiver latency, raw loopback excluding logic and protocol \\
Exegy (incl.\ Enyx, acquired 2022)~\citep{exegy_enyx,exegy_stac2024} & FPGA feed handling, order entry and trigger frameworks & STAC-T0 audited trigger latency of the order of ten nanoseconds \\
Algo-Logic~\citep{algologic} & CME tick-to-trade system on Alveo U50, with feed handler, book and TCP order entry & under 500\,ns tick-to-trade \\
Magmio~\citep{magmio} & C++ strategies compiled to FPGA by high-level synthesis & under 100\,ns for triggers; under 350\,ns with book building \\
Xelera Silva~\citep{xelera_silva} & tree-ensemble and small-network inference on FPGA & about 250\,ns (99th percentile) for a gradient-boosted model inline on UL3524; 1.1--1.3\,\textmu s as card offload; audited STAC-ML: at most 1.95\,\textmu s (Table~\ref{tab:stac}) \\
Napatech + Xelera~\citep{napatech} & the same inference on FPGA network cards & 1.55\,\textmu s for a 1000-tree model with 100 features \\
\bottomrule
\end{tabularx}
\end{table}

\subsection{Accelerators between FPGAs and GPUs}

A class of accelerators aims at the gap between FPGA determinism and GPU flexibility. This section gives three
examples.
LightTrader~\citep{lighttrader}, from KAIST and the chip designer Rebellions, combines an FPGA for
market data and orders with dedicated AI processors and a scheduler that manages workload and power;
back-tested on CME data, it reports speed-ups of about 14 times over a GPU-based and 7 times over an
FPGA-based system for AI inference. Corsair~\citep{corsair2025} is an in-memory-computing design for large-language-model inference, built from
chiplets (small dies combined in one package); it is not a trading product, but it illustrates the
direction of the hardware industry. Network cards with programmable logic carry inference into the
network path itself~\citep{napatech}.

\subsection{Transformers on FPGAs}

Mapping attention to FPGAs is an active research topic outside finance; this section summarises two designs and
what they imply for order-book Transformers. FAMOUS~\citep{famous}
accelerates dense multi-head attention on a data-centre FPGA, at 0.6 to 0.94\,ms per attention layer for sequences
of 64 elements; its comparison with a V100 GPU uses a figure taken from another paper rather than measured.
ELiTeFormer~\citep{eliteformer} co-designs a ternary-weight Transformer with linear attention for an adaptive system-on-chip (a chip that combines processor cores with programmable logic).
Neither is a trading system, and we found no measured FPGA implementation of an order-book Transformer. Small,
quantised attention models are candidates for programmable logic, with quantisation as the enabling
technique~\citep{jacob2018}; full spatial-temporal Transformers remain GPU workloads.

\subsection{GPUs for training, simulation and single-event inference}

Graphics processors are built for data-parallel workloads: training, large-scale simulation, and models
whose memory footprint or operator mix exceeds practical FPGA budgets. JAX-LOB~\citep{frey2023}
(Section~\ref{sec:sim-survey}) moves order-book matching onto accelerators for reinforcement-learning training. Monte Carlo simulation
of order-book dynamics parallelises well across paths; one vendor demonstration reports speed-ups of
14 to 114 times over a CPU~\citep{nvidia_numba2025}, though for simulated paths rather than a
historical matching-engine backtest. Specialised processors also accelerate training: a multi-horizon
DeepLOB variant trained in about 15\% of the GPU time on intelligence processing
units~\citep{zhangzohren2021}. For single-event inference, GPUs used naively are slow (Section~\ref{sec:hw-why}),
but with vendor-optimised code and a sliding window that reuses earlier computation, a GH200 reached 4.70\,\textmu s
at the 99th percentile on the smallest audited LSTM (Table~\ref{tab:stac}).

\subsection{Networks and distance}

Latency is not only computation. This section covers the time spent in the operating system's network stack, in
the fibre, and in network devices that compute. A peer-reviewed measurement of the latency budget of a server gives medians of
4.5\,\textmu s through the operating system's network stack, 946\,ns with kernel bypass, and 572\,ns for one small
read across PCIe~\citep{zilberman2017where}. The operating system's network stack adds microseconds;
kernel-bypass libraries deliver packets directly to the application~\citep{dpdk}. Published figures for
such stacks include half-round-trip times of about 3\,\textmu s on older 10-gigabit
hardware~\citep{openonload} and about 1\,\textmu s for UDP~\citep{vma}. Physical distance sets a floor no
software can remove: light in optical fibre takes about 5\,\textmu s per kilometre (about 4.9\,ns per
metre)~\citep{corning_smf28,fiberlatency}, which is why trading systems are colocated in the
exchange's data centre and why differences in relative speed translate into differences in trading
performance~\citep{baron2019,budish2015}.

Computation can also move into the network itself. Programmable switches have filtered a Nasdaq ITCH feed by
symbol, delivering all messages of a recorded trace within 50\,\textmu s against 300\,\textmu s for the software
baseline~\citep{jepsen2018packet,jepsen2020forwarding}, and machine-learning models have been run inside switches on
ITCH data~\citep{zheng2024planter}. We found no measured order book or matching engine built inside a switch. At the extreme, application-specific chips fix a design in
silicon, trading the reprogrammability of an FPGA for speed and power~\citep{kuon2007}.

\subsection{Placing workloads on hardware}

Table~\ref{tab:placement} gathers the preceding subsections into a qualitative guide to which hardware each
workload belongs on.

\begin{table}[!htbp]
\centering\small
\caption{A qualitative guide to workload placement. It reflects the general division of labour
between platforms, not measured results.}
\label{tab:placement}
\begin{tabularx}{\linewidth}{lXX}
\toprule
Workload & Preferred platform & Rationale \\
\midrule
Market-data parsing, book maintenance & FPGA & regular, protocol-bound, deterministic \\
Linear OFI, queue imbalance, simple triggers & FPGA or CPU & few operations; easy to pipeline \\
Small quantised networks or tree ensembles on the critical path & FPGA & feasible when capacity is
  constrained by design \\
Pre-trade risk checks, mass cancel & FPGA & must be deterministic and next to order generation \\
Multi-horizon or cross-sectional deep models, full Transformers & GPU or other accelerator &
  memory and operator mix exceed FPGA budgets \\
Large-scale simulation, training, hyper-parameter search & GPU & throughput dominates \\
Strategy logic that changes often & CPU with kernel bypass & flexibility outweighs nanoseconds \\
\bottomrule
\end{tabularx}
\end{table}

The guiding rule is Occam's razor (Principle~2) again: place a computation on the most constrained, and
therefore most deterministic, hardware that still delivers the required economic performance.
Moving a model to faster hardware is justified only when the reduction in latency and jitter
measurably improves post-cost performance.

\subsection{Hardware cannot extend signal lifetime}

Simpler models have a structural advantage on deterministic hardware: fewer parameters and
simpler data flow map onto fixed-point, branch-free pipelines. Acceleration can rescue part of a
complex model's value by removing engineering latency, but it cannot extend the statistical
lifetime of the signal.

\begin{example}[Second-order latency]
A logistic score on five OFI windows and queue imbalance takes a few tens of multiply--add
operations. If the signal's half-life is 5 ms, evaluating it in 50\,$\mu$s rather than 500\,ns
keeps $2^{-0.01} \approx 99.3\%$ of its value rather than essentially all of it. The difference is
second-order; the binding constraint is the quality and horizon of the signal itself.
\end{example}

The design sequence that follows is: identify the horizon and value of the signal under idealised
execution; choose the simplest model that captures that value; only then choose the hardware tier
whose latency and jitter preserve it after costs.

% ================================================================
\section{Summary of Part III}
\label{sec:part3-summary}
% ================================================================

Part~III turned from models to the infrastructure needed to test and run them. This section collects its main points;
it adds no new results.

\paragraph{What a simulator must capture.} A passive order earns its average profit as the product of two factors:
the probability that it is filled and its average profit when it is (Equation~\eqref{eq:decomp}). Both depend on
queue position. Fills are biased towards bad ones, because an order is most likely to be filled when an informed
trader wants to trade against it. Latency eats into a signal: with a 2\,ms half-life, 97\% of its value remains after
100\,\textmu s and 18\% after 5\,ms (Section~\ref{sec:sim-needs}).

\paragraph{Simulators.} A backtest that fills an order when the price touches it fills orders the market would never
have reached. Simulators differ in where the market comes from---replay, agents, a fitted model or a generative
network---and in how they place the simulated order in the queue. Agent-based, model-based and generative simulators
let the market react to the orders, but their queues are only as realistic as the simulated order flow. Among replay
tools, three place the simulated order in the recorded queue without changing the recorded market: hftbacktest, lobsim
and Kaspar. The commercial tools whose documentation we could read estimate queue position, fill without a queue, or
do not say how they do it (Section~\ref{sec:sim-survey}).

\paragraph{Kaspar.} Kaspar replays an exchange's own market-by-order feed (CME MDP~3.0 today), places simulated orders
in the reconstructed first-in-first-out queues after separate order, cancel and feed delays, fills them only when
recorded executions reach an order behind them, and runs the same code in replay, paper and live trading. It replays
one day per process, so many days can run at once on one machine; among commercial replay tools only Sierra Chart's
licence allows the same (Table~\ref{tab:parallel}). A model of Part~II is added as one more actor that reads the book
and sends its predictions to the strategy (Section~\ref{sec:kaspar-actor}).

\paragraph{Hardware.} Book updates take about 8\,\textmu s in Kaspar's software path and a few hundred nanoseconds or
less on FPGAs; audited FPGA tick-to-trade paths with simple logic run in about 0.1 to 0.6\,\textmu s. Published
inference times of order-book models range from about 25\,\textmu s to 23\,ms (Table~\ref{tab:lob-latency}). Four
causes make models slow on a single event: recurrence that runs step by step, attention whose cost grows with the
square of the window, memory traffic at a batch of one, and fixed overheads per GPU operation
(Section~\ref{sec:hw-why}). With a sliding window that reuses earlier computation, audited inference of a small LSTM
takes single-digit microseconds on both FPGAs and GPUs (Table~\ref{tab:stac}). No audited benchmark yet evaluates a
learned model on a book it has built.

% ================================================================
\section{Research agenda}
\label{sec:agenda}

The survey has identified places where the literature stops short of what a trader needs. They are grouped below
by the part of the survey they come from.

\subsection{Models}

The open problems in this group concern what the models of Part~II predict and how they are built.

\paragraph{Execution-aware objectives.} Most prediction models stop at
an estimate of $\P(\Delta m_{t,H} > 0 \mid X_t)$. The more useful objective is an estimate of
$\E[\Pi \mid a_t, X_t]$, the expected profit of each available action,
combining prediction, fill probability, queue position and adverse selection. Calibrated
confidence~\citep{manoharan2026uqlob} is a step toward it: it tells an execution policy when to act.

\paragraph{Message-level foundation models for execution.} Pretrained message-level encoders have
so far been evaluated mainly on mid-price and next-message prediction~\citep{linna2025lobert}.
Fine-tuning them for fill probability and conditional mark-outs, evaluated in queue-exact replay,
is an open direction.

\paragraph{State-space and long-context models for event streams.} State-space sequence models are well
suited to long, irregular event streams, yet apart from one hardware-oriented design~\citep{popov2026ssm} we found no established LOB forecaster built on them.
Testing them against simple baselines on the same inputs---not against each other---is a clear
next step.

\paragraph{Near-criticality and learned models.} Estimated branching ratios of order flow are high,
from about $0.7$--$0.8$ to about one depending on method. How deep sequence
models behave on near-critical, long-memory inputs, and whether Hawkes structure built into them
(Section~\ref{sec:hybrid}) improves robustness, is untested.

\paragraph{A parsimony audit.} For each widely cited deep LOB model, re-estimate its incremental
economic value and deflated performance against the best low-capacity baseline on the same data and
fill model. Such a meta-study would directly test the central claim of this survey.

\subsection{Execution and simulation}

The open problems in this group concern how predictions become orders and how simulators measure the result
(Sections~\ref{sec:sim-needs} and~\ref{sec:simulators}).

\paragraph{Queue-aware benchmarks.} Queue position should be a first-class state variable.
Benchmarks should report performance conditional on queue rank, residual volume, spread and expected
depletion time, and should measure the size of the queue-rank effect on fill probability and on
conditional mark-outs across assets and regimes, rather than take a single model estimate as
given.

\paragraph{End-to-end signal and execution.} Most models output a forecast and leave execution to a
separate rule. Models whose output is the decision itself---quote, keep, cancel, cross---and whose
training objective is execution value would close the gap that this survey describes. Reinforcement
learning in LOB simulators~\citep{lalor2025nhp,frey2023} is one route, ideally validated in queue-exact replay; its results must be
judged against model-free benchmarks in the same simulator.

\paragraph{Causal rather than predictive questions.} The questions that matter for execution are
causal: they ask what happens to the queue and the price \emph{if} an order is placed at a given price. Replay with a
no-impact assumption answers this only for small orders. Models of the market's response to one's own
orders, and ways to validate them without live experiments, remain open.

\paragraph{Generative evaluation.} Generators should be evaluated by market invariants and downstream
utility rather than by visual or marginal similarity alone~\citep{lobbench2025}.

\paragraph{Cross-asset information after costs.} An open question is when cross-asset information remains
valuable after costs, latency, multiple testing and parameter instability.

\paragraph{Simulators against live fills.} The simulators of Table~\ref{tab:simulators} have not been compared on the
same orders. Running one strategy through several of them and through paper trading, and comparing fill rates and
post-fill mark-outs with the live fills, would measure how much each simulator's assumptions---estimated queue
position, simulated order flow, no market impact---cost in accuracy.

\paragraph{Replay with a reacting market.} Replay cannot measure the market's response to one's own orders; a
model-based simulator can, mechanically (Section~\ref{sec:kaspar-fills}). Combining the two---recorded flow for the
rest of the market, a fitted model for the response to the strategy's orders---and validating the response against
live trading is open.

\subsection{Hardware}

The open problems in this group follow from Section~\ref{sec:hardware}.

\begin{itemize}[leftmargin=1.5em,itemsep=1pt]
\item Training objectives that penalise expected inference latency and resource use on a target
  device, so that models trade accuracy against deployability explicitly.
\item Which components of spatial-temporal LOB models survive mapping to programmable logic without
  losing their economic value.
\item Hawkes intensities in hardware: exponential kernels update in constant time and suit pipelines,
  and a power-law kernel can be approximated by a few exponentials, each with its own running
  sum~\citep{bochud2007}; how many are needed at the precision of fixed-point hardware, and how
  state-dependent kernels map to pipelines, are open.
\item Whether the statistical properties of a strategy developed on one hardware and latency profile
  survive deployment on another---the research--production gap of Section~\ref{sec:kaspar-determinism}, at the hardware layer.
\item Public, auditable benchmarks that cover book building and inference together: STAC-T1 audits a whole
  tick-to-trade path with simple logic, and STAC-ML audits inference alone, excluding data ingestion and features;
  neither evaluates a learned model on a book it has built.
\end{itemize}

% ================================================================
\section{Conclusion}
\label{sec:conclusion}
% ================================================================

This section summarises the survey's argument, the two kinds of infrastructure needed to test it, and the order of
operations it recommends to a reader who must choose a model.

The literature on limit-order-book modelling has moved from low-dimensional queue and order-flow
models to expressive sequence architectures and pretrained message-level encoders. Statistical
prediction has improved. Statistical prediction and executable trading remain different problems.
The translation between them is
\[
  \text{LOB state} \to \text{representation} \to \text{prediction} \to \text{decision} \to
  \text{queue} \to \text{execution} \to \text{P\&L}.
\]
\inwords{the book is turned into model inputs, the inputs into a forecast, the forecast into an order, the order
waits in a queue and may be filled, and the fill produces a profit or loss. Each link can destroy information or add
uncertainty.}

Order-flow imbalance, queue imbalance, Hawkes intensities and first-passage models give
interpretable low-dimensional descriptions of market dynamics. Deep and pretrained models give
representational capacity and, increasingly, calibrated confidence; most of them are learned analogues of
the classical signal models rather than of the mechanistic ones (Section~\ref{sec:deep-adds}). Meta-models and
structural hybrids combine the two, either by feeding classical outputs to a learner or by building classical
structure inside a network. Much of the accuracy reported for deep models depends on how the label and the data
split are built: with a label that averages past prices, a re-test of DeepLOB on Apple was right 65.9\% of the time
and 54.6\% with a label measured from the current price~\citep{lee2024predictability}, and in our synthetic example a
network reached 57.6\% on a market with nothing to predict (Section~\ref{sec:deep-train}). The choice of input
matters as much: with the network fixed, order-flow inputs were among the best in 89\% of cases and raw snapshots in
11\%~\citep{lucchese2024}. The conclusion is not that simple models are superior. It is that complexity is an
empirical hypothesis: a more complex model is justified when it delivers incremental out-of-sample
economic value that survives realistic execution, temporal validation, and the multiple-testing
burden of its own selection.

Two kinds of infrastructure make that standard testable. The first is the simulator. Agent-based, model-based and
generative simulators let the market react to the researcher's orders, but their queues are only as realistic as the
simulated order flow; a replay keeps the real market but cannot react. Exact queue position in the recorded market
needs order-level data and an order-level book; among the replay tools compared, hftbacktest (with order-level
data), lobsim and Kaspar provide it without altering the recorded market (Section~\ref{sec:sim-survey}). In a replay built from actors, a model of
Part~II becomes one more actor that reads the book and sends its predictions to the strategy, and the same code runs
in replay, paper and live trading (Section~\ref{sec:kaspar-actor}).

The second is hardware. Book building takes about 8\,\textmu s from socket to book in Kaspar's software path and
hundreds of nanoseconds or less in published FPGA designs; published inference times of order-book models range from
about 25\,\textmu s to about 23\,ms, about half of them above one millisecond (Table~\ref{tab:lob-latency}). Section~\ref{sec:hw-why}
identifies four causes---recurrence that must run step by step, attention whose cost grows with the square of the
window, memory traffic at a batch of one, and fixed overheads per operation---and Section~\ref{sec:hw-what} lists the
candidate remedies: shorter or incremental windows, smaller and lower-precision models, fused kernels, and moving the
regular stages to FPGAs. Hardware
decides how much of a signal's value survives to the moment of execution, but it cannot rescue a signal whose life is
shorter than the pipeline that computes it.

For a reader who must choose a model, the survey's guidance reduces to an order of operations.
Decide first what the output is for---aggressive trading, passive execution, market making or
simulation---and therefore which of the targets of Section~\ref{sec:deep} matters. Establish the
strongest simple baseline on the best available representation: order-flow imbalance, queue
imbalance, the micro-price, the queue race, and for passive orders the volume ahead. Evaluate it with
proper scores and economic metrics in a queue-exact replay, counting every configuration tried. Only
then add capacity---a shallow network, a sequence model, a hybrid---and keep it only if the increment
survives the same test.

A useful future literature should therefore report not only whether a model predicts the market,
but what information it extracts, what decision it changes, how that decision is executed, and
whether the incremental value survives a properly constructed null.

\begin{table}[!htbp]
\centering\small
\caption{Evidence status of the main claims in this survey.}
\label{tab:evidence}
\begin{tabularx}{\linewidth}{Xl}
\toprule
Claim & Status \\
\midrule
Short-horizon mid changes are approximately linear in OFI, with slope inversely related to
  depth~\citep{cks2014} & established \\
Hawkes stationarity when $\rho(\Gamma) < 1$ and $\bar\lambda = (I-\Gamma)^{-1}\mu$ & established \\
Input representation can matter as much as the architecture & synthesis \\
Naive ``touch'' backtests overstate passive fills relative to queue-exact replay & synthesis;
  argued from price--time priority \\
Front-of-queue positions are worth more than back-of-queue positions, with adverse selection
  increasing in queue position~\citep{moallemi2017} & established in a model; size across regimes open \\
Hybrid inputs derived from the same stream add little to a sufficiently large network & open
  hypothesis \\
Labels that average past prices inflate reported direction accuracy~\citep{lee2024predictability} & one re-test on
  real data; reproduced on synthetic data here \\
Order-flow inputs beat raw snapshots with the network held fixed~\citep{lucchese2024} & established in one study
  (10 stocks) \\
Published inference times of LOB models range from tens of microseconds to tens of milliseconds, above book-update times & measured in five
  studies under different conditions \\
Exact own-order queue position in the recorded market requires order-level data and an order-level book &
  synthesis; argued from price--time priority \\
\bottomrule
\end{tabularx}
\end{table}
